\section{Background}

\subsection{Safety Cases and Black-Box Logs}
Static safety cases~\cite{kelly1998,easa_sc_vtol} argue that a
system is safe by design. Black-box logs (ArduPilot~\cite{ardupilot},
PX4~\cite{px4}) record what actually happened during a mission:
trajectory, sensor readings, commands. Neither answers the
counterfactual question: \emph{what would have happened if conditions
had been different?}

\subsection{Counterfactual Reasoning}
Counterfactual reasoning~\cite{pearl2009,lewis1973} is the formal
study of ``what if?' questions. Given a causal model and an
observation, it computes the outcome under a hypothetical
intervention. Wachter et al.~\cite{wachter2017} apply
counterfactuals to explain black-box models, but not to certify
them. CSR applies counterfactuals to the \emph{mission}, not to a
model: the causal model is the physics, the observation is the
recorded trajectory.

\subsection{Cryptographic Audit}
A hash chain~\cite{crosby2009,merkle1988} makes a log
tamper-evident. CSR extends the chain with counterfactual proofs:
each proof is bound to the chain root, so that tampering with the
chain invalidates the proof, and vice versa.

\section{CSR Design}

\subsection{Trajectory Replay}
CSR takes the recorded trajectory (a sequence of
\texttt{TrajectorySample} structs with time, position, velocity,
wind, and altitude) and replays it under modified conditions. Two
modifications are supported:
\begin{itemize}
    \item \texttt{what\_if\_wind(scale)}: multiply the wind vector
          by a scalar. The drift is integrated over the trajectory.
    \item \texttt{what\_if\_altitude(offset)}: shift the altitude
          by an offset. The outcome is classified as
          \texttt{Failed} if the result is outside the envelope
          ($<0.5$~m or $>120$~m).
\end{itemize}
Each result is classified as \texttt{SameAsActual}, \texttt{Better},
\texttt{Worse}, or \texttt{Failed}.

\subsection{Proof Construction}
A counterfactual result is bound to the original chain by:
\begin{verbatim}
result_hash = SHA-256(scenario || modification || parameter
                      || waypoint_error)
proof_hash  = SHA-256(chain_root || result_hash)
\end{verbatim}
Verification recomputes \texttt{result\_hash} and \texttt{proof\_hash}
and compares them to the stored values. If the chain has been
modified, \texttt{chain\_root} changes, and the proof fails.

\section{Implementation}

\subsection{Crate Structure}
CSR is implemented as the crate \texttt{rnfc-counterfactual} (250
lines of Rust). The public API has three functions:
\begin{itemize}
    \item \texttt{what\_if\_wind(traj, target, scale) -> CounterfactualResult}
    \item \texttt{what\_if\_altitude(traj, offset) -> CounterfactualResult}
    \item \texttt{prove(chain, result) -> CounterfactualProof}
\end{itemize}
The crate depends on \texttt{rnfc-certchain} for the chain and on
\texttt{sha2} for the hash.

\subsection{Integration with CertChain}
The original mission produces a \texttt{CertChain} with one block per
safety event. CSR appends a new block for each counterfactual proof,
so that the chain contains both the actual mission and the
counterfactual analysis. A regulator can verify the whole chain in
one pass.

\section{Evaluation}

\subsection{Case Study: 4 Scenarios}
We evaluate CSR on a mission with target $[50, 0]$ and actual final
position $[50, 1.5]$ (error 1.5~m). Four scenarios are analysed:
\begin{itemize}
    \item \texttt{wind\_0x}: no wind. Outcome \texttt{Better}
          (error 0.0~m).
    \item \texttt{wind\_1x}: baseline. Outcome \texttt{SameAsActual}
          (error 1.5~m).
    \item \texttt{wind\_20x}: extreme wind. Outcome \texttt{Failed}
          (error 30.0~m, above the 10.0~m limit).
    \item \texttt{altitude\_+150m}: altitude +150~m. Outcome
          \texttt{Failed} (altitude 160~m, above the 120~m EASA
          limit).
\end{itemize}

\subsection{Security Properties}
Three attacks were tested and all rejected:
\begin{itemize}
    \item \textbf{Tampered chain}: modifying a block changes
          \texttt{chain\_root}, so \texttt{proof\_hash} no longer
          matches. Verdict: \texttt{BadSignature}.
    \item \textbf{Tampered result}: modifying \texttt{waypoint\_error}
          changes \texttt{result\_hash}, so the proof fails.
          Verdict: \texttt{BadSignature}.
    \item \textbf{Wrong key}: using a different key changes
          \texttt{proof\_hash}. Verdict: \texttt{BadSignature}.
\end{itemize}

\subsection{Limitations}
The physics is synthetic (1D/2D), not 6-DOF. The hash is FNV-1a
truncated to 64 bits in the current implementation, not full
SHA-256 (the proof construction uses SHA-256, but the chain itself
uses FNV-1a). There are no formal guarantees. CSR is a proof of
concept.

\section{Discussion}

\subsection{Relation to DO-326A}
CSR provides evidence for the post-mission audit requirement
(SEC-12 in the project's security matrix). A regulator can ask
``what if the wind had been $2\times$?' and receive a proof that
the answer is correct, without trusting the operator. This
complements the static safety case and the runtime monitors.

\subsection{Ethical and Legal Considerations}
Who owns the counterfactual? Who can request it? If a drone
crashes, and the counterfactual shows that a different route would
have been safer, is the operator liable? These questions are
outside the scope of this paper, but they must be addressed before
CSR is used in certification. We note them as open problems.

\section{Conclusion}
We presented CSR, a technique that enables a drone to answer
post-mission ``what if?' questions with cryptographically
verifiable answers. CSR replays the mission under modified
conditions and binds the result to the original chain via SHA-256.
On a 4-scenario case study, 4/4 proofs verified and 3/3 attacks
were rejected. CSR is a proof of concept; future work includes
real 6-DOF physics, full SHA-256 in the chain, and formal
guarantees.
